\chapter*{ABSTRACT}
\addcontentsline{toc}{chapter}{ABSTRACT}

People with mild cognitive impairment (MCI) vary in how quickly they experience further cognitive or functional decline. To rank participants by progression risk, this study combines magnetic resonance imaging (MRI) at three time points with quantitative features extracted from the scans (radiomics) and clinical scores.

At each visit, the 0 CP model encodes MRI, radiomics and clinical data separately, normalizes the three representations, concatenates them and projects them into a shared visit representation. The name 0 CP indicates that all three explicit pairwise interactions are disabled. A time-aware LSTM processes the three visits, attention pools their hidden states, and a Cox head produces a risk score.

Eight fusion architectures were compared on the same split of 359 ADNI participants (251 training, 54 validation and 54 test). Across three initializations, 0 CP had the highest mean validation C-index ($0.8434\pm0.0288$) and was selected as the primary model. RC-Free, which retains both MRI-based interactions, achieved $0.8347\pm0.0237$. After selection, 0 CP achieved a C-index of $0.8914\pm0.0097$ on the previously accessed test set.

\noindent\textbf{Keywords:} Alzheimer's disease, MCI, longitudinal MRI, radiomics, pairwise fusion, survival analysis.
